% Compile with: uplatex file.tex && dvipdfmx file.dvi
%           or: latexmk -uplatex file.tex
%\documentclass[uplatex,dvipdfmx]{jsarticle}
\documentclass[dvipdfmx,a4paper,11pt]{jsarticle}
\usepackage{amsmath, amssymb}
\usepackage{tikz}

\title{scenario\_free\_electron\_fermi\_surface}
\date{\today}

\begin{document}
\maketitle

フリー電子モデルの分散関係を扱う。波数 $k$、電子質量 $m$、化学ポテンシャル $\mu$、換算プランク定数 $\hbar$ をスカラーパラメータとして準備する。

k, m, $\mu$, $\hbar$ をスカラーとして定義する

\[
  k, m, \mu, \hbar
\]

フリー電子のエネルギー分散 $\xi_k=\hbar^2k^2/(2m)-\mu$（化学ポテンシャルを基準としたエネルギー）を定義する。

$\xi_k = \hbar^2 k^2 / (2 m) - \mu$

\begin{equation}
  \xi_k = \hbar^2 k^2 / (2 m) - \mu
  \label{eq:E1}
\end{equation}

右辺を整理して見やすい形にする（物理的な意味は変わらない）。

E1 の右辺を整理して

\begin{equation}
  \xi_k = \frac{ - 2 m \mu + k^{2} \hbar^{2}}{2 m}
  \label{eq:E2}
\end{equation}

フェルミ面の条件 $\xi_k=0$（エネルギーがフェルミ準位に一致する）を、右辺を 0 と置いた方程式として確定する。

E1 の右辺を 0 とする

\begin{equation}
  \hbar^2 k^2 / (2 m) - \mu = 0
  \label{eq:E3}
\end{equation}

この方程式を波数 $k$ について解き、フェルミ波数 $k_F$ を求める。

E3 を k について解いて

\begin{equation}
  k = \sqrt{2} \sqrt{m \mu}/\hbar
  \label{eq:E4}
\end{equation}

得られたフェルミ波数を元の $\xi_k$ の式に代入し、フェルミ面上ではエネルギーが恒等的にゼロになることを確認する。

E1 に E4 を代入して

\begin{equation}
  \xi_k = \left( - \mu + \frac{\left( \frac{\sqrt{2 m \mu}}{\hbar} \right)^{2} \hbar^{2}}{2 m}\right)
  \label{eq:E5}
\end{equation}

代入結果を整理し、ゼロになることを明示的に確認する。

E5 の右辺を整理して

\begin{equation}
  \xi_k = 0
  \label{eq:E6}
\end{equation}

有効質量 $m^*$ を求めるため、分散関係 $\xi_k$ を波数 $k$ について2階微分する。

E1 を k で 2 回微分して

\begin{equation}
  \frac{d^{2}}{dk^{2}}\left(\xi_k\right) = \frac{\hbar^{2}}{m}
  \label{eq:E7}
\end{equation}

得られた2階微分から、有効質量 $m^*=\hbar^2/(d^2\xi_k/dk^2)$ を定義する（フリー電子モデルでは $m^*=m$、裸の質量に一致するはず）。

$m_star = \hbar^2 / E7$

\begin{equation}
  m_star = m
  \label{eq:E8}
\end{equation}

\end{document}